\begin{pastexam}[label={px:ntu-109-comm-dsa-44}]{臺大電信所 109 學年度（B）第 44 題}
（複選；每題 2 分，每答錯一個選項倒扣 0.4 分）

In a splay tree of size n
\begin{enumerate}[label=(\Alph*),nosep]
\item The height is $O(\log n)$.
\item Time complexity of one insertion is $O(\log n)$.
\item Time complexity of n insertions is $O(n\log n)$.
\item Time complexity of one deletion is $O(\log n)$.
\item Time complexity of n deletions is $O(n\log n)$.
\end{enumerate}
\end{pastexam}

\begin{pastsoln}
答案：(C)(E)。

伸展樹不保證高度，只保證\textbf{平攤}成本：每次操作平攤 $O(\log n)$（資結書定理~10.6）。
單一操作的最壞情形可以很差，一串操作的總和才有保證。

\textbf{(A) 錯。} 依序插入 $1, 2, \ldots, n$：每次新的鍵比樹裡所有鍵都大，插在根的右邊，
再做一次 zig 轉到根，舊的根變成它的左子節點。 最後的樹是 $n$ 往左一路到 $1$ 的一條鏈，高度
$n - 1$。

\textbf{(B) 錯。} 接著在上面那棵鏈狀的樹插入 $0$，要從根一路往左走到 $1$ 底下，走 $n$ 步，再伸展
$n$ 層，單次是 $\Theta(n)$。

\textbf{(C) 對。} 設位能 $\Phi(T) = \sum_x \lg s(x)$（$s(x)$ 是子樹大小）。 $n$ 次插入的實際總
成本 $= \sum(\text{平攤成本}) - \Phi(T_{\text{末}}) + \Phi(T_0) \le \sum(\text{平攤成本}) + \Phi(T_0)$。
過程中樹最多 $2n$ 個節點，每次平攤 $O(\log 2n) = O(\log n)$；一開始的樹有 $n$ 個節點，
$\Phi(T_0) \le n \lg n$。 合計 $O(n \log n)$。 從空樹開始也一樣。

\textbf{(D) 錯。} 和 (B) 同理：刪除要先搜尋，在鏈狀的樹裡刪最深的節點，單次 $\Theta(n)$。

\textbf{(E) 對。} 同 (C)：每次刪除平攤 $O(\log n)$，$n$ 次共 $O(n\log n)$，再加上
$\Phi(T_0) \le n\lg n$，仍是 $O(n\log n)$。
\end{pastsoln}
